\section{Introduction}
\label{sec:introduction}

Solubility is an intrinsic characteristic of a compound, usually defined as maximal achievable equilibrium concentration of a compound (solute) in a certain solvent under given conditions (temperature, pressure). The knowledge of solubility values for a single solute-solvent pair dovetails with its further applications in drug design [\cite{ma2025solecos,bolla2022crystals}], organic synthesis pipelines [\cite{reichardt2011solvents,diorazio2016toward,tu2025askcos}] or crystallization studies [\cite{sheikholeslamzadeh2012optimal,mendis2022simultaneous}]. Namely, water solubility usually guides the preparation of pharmaceutically active substances, whereas organic solubility determines solvent choice in flow and batch synthesis, affecting entire sectors of the economy.

Thus, any \textit{a priori} knowledge on solubility of organic compounds in organic solvents and water will notably accelerate the development of the abovementioned technological processes. For this purpose, many predictive approaches were developed so far to estimate solubilities across diverse subsets of compounds [\cite{fowles2025physics}], achieving several noticeable milestones. Boobier \textit{et al.} incorporated a set of computational physicochemical descriptors, which enabled creation of simple yet performative and predictive machine learning models [\cite{boobier2020machine}]. Later, Attia \textit{et al.}[\cite{attia2025data}] and Al Ibrahim \textit{et al.}[\cite{al2025accurately}] developed neural network architectures which performance is approaching the so-called aleatoric limit (inherent error caused by inconsistencies in the data). Numerous other approaches incorporating large language models, graph transformers or even boosting algorithms with carefully distilled set of features can reach close to this limit [\cite{jung2025enhancing,broadbent2025solvaformer,ramani2026dissolvr}]. It should be highlighted that most of these methods rely on large-scale experimental datasets, with AqSolDB [\cite{sorkun2019aqsoldb}], BigSolDB2.0 [\cite{krasnov2025bigsoldb}] and MixtureSolDB  [\cite{malikov2026dataset}] being largest congeners up to date.

Still, although notable success was achieved in terms of models performance, severe discrepancies in data collection, cleaning and sanitation as well as in validation strategies across different studies hinder direct comparison of methodologies, what hampers fair assessment of the realm of developed approaches. Moreover, imprecise data handling can limit our insights in terms of actual generalization capability of ML models, their intrinsic interpretability and real-life application potential [\cite{llompart2024curation}]. To this end, a thorough benchmark study based on a large-scale experimental database is needed; \textsc{BigSolDB}, currently the largest literature-aggregated solubility resource, provides that foundation.

In this work we introduce \scthree{}, an open multi-solvent benchmark derived from \textsc{BigSolDB}~v2.1 after systematic auditing and cleaning (§\ref{sec:data_curation}), with an empirically grounded aleatoric floor and tiered consensus evaluation sets carrying per-point uncertainties (§\ref{sec:aleatoric}), leakage-checked train/eval/OOD/tier splits (§\ref{sec:benchmark}), and a metric suite tailored to multi-solvent regression (§\ref{sec:metrics}). We train and evaluate a broad catalogue of baselines under fixed protocols (§\ref{sec:baselines}--\ref{sec:results}), analyse data scaling relative to the calibrated floor (§\ref{sec:data_scaling}), interpret models where labels allow meaningful attribution (§\ref{sec:interpretability}), and study transfer from an adjacent quantum-chemistry solvation task (§\ref{sec:transfer}). Source code, split files, and regeneration scripts accompany the paper.
